% ============================================================
%  封面信息
% ============================================================
\school{计算机科学与技术学院}
\major{软件工程}
\student{2250758}{林继申}
\thesistitle{基于AI Agent的建筑幕墙韧性评估软件系统设计与实现}{}
\thesistitleeng{Design and Implementation of an AI Agent-Based Software System for Building Curtain Wall Resilience Assessment}{}
\thesisadvisor{黄杰}
\thesisdate{2026}{5}{1}

% ============================================================
%  信息说明页
% ============================================================

% 成果类型：design（毕业设计）或 thesis（毕业论文）
\infotype{design}

% 内容简述（请用300字以内简要概述）
\infoabstract{%
本毕业设计面向建筑幕墙韧性评估中人工巡检效率低、结果主观、图像数据难以规模化处理和报告难以系统生成的问题，设计并实现一套基于 AI Agent 的软件系统。系统以项目级图像资产管理为基础，由 Agent 自主判断每张图像需要执行的原子检测能力，并通过分层信息整合策略先生成单图总结，再融合为项目级评估报告。工程上采用前后端分离、gRPC、RocketMQ、WebSocket、MySQL、Redis 与 MinIO 等技术，形成可部署、可追踪、可扩展的端到端原型。本文重点在于 Agent 赋能工程链路，而非单个视觉模型本身。
}

% 毕业设计：图纸数量与正文字数（成果类型为 design 时填写，两者须同时填写）
\infodrawings{5}
\infowordcount{27730}

% 毕业论文：正文总字数（成果类型为 thesis 时填写）
\infothesiswords{12345}

% 随附资料（每条用 \item 开头；不填则显示空白行）
\infomaterials{
  \item 建筑幕墙韧性评估软件系统源代码；
  \item 数据库定义、接口定义与部署配置文件；
  \item 图像检测样例、运行说明与系统演示材料。
}
